\documentclass[]{exam}
\usepackage{amssymb}
\usepackage[top=1in, bottom=1.5in, left=1in, right=1in]{geometry}
\usepackage{tikz}
\usepackage{pgfplots}
\checkboxchar{$\Box$}
\usepackage{graphicx}
\usepackage{caption}
\usepackage{float}
\usepackage{enumerate}
\usepackage{comment}
\usepackage{multicol}
\usepackage{color}
\CorrectChoiceEmphasis{\color{red}}
\SolutionEmphasis{\color{red}}
\usepackage{advdate}
\usepackage{gensymb}
\usepackage{mathrsfs}
\usetikzlibrary{arrows}


\usepackage{pgfkeys,pgfcalendar}

\newcount\pgfdatecount
\newcommand{\tomorrow}{%
\pgfcalendardatetojulian{\year-\month-\day+1}{\pgfdatecount}
\pgfcalendarjuliantodate{\the\pgfdatecount}{\myyear}{\mymonth}{\myday}
\pgfcalendarmonthname{\mymonth}\space\myday,\space\myyear%
}


\newcommand{\advanceday}[1][1]{
\begingroup
\AdvanceDate[#1]
\today
\endgroup}

\newcommand{\vertex}{\node[vertex]}
\tikzstyle{vertex}=[circle, draw, inner sep=0pt, minimum size=6pt]

\usepackage{lastpage} 
\usepackage{extramarks}
\usepackage{graphicx} 
\usepackage{tikz}
\usepackage{sectsty} 
\usepackage{amsmath,amsfonts,amsthm} 
\usepackage{wrapfig, blindtext}

% Margins
\topmargin=-0.45in
\evensidemargin=0in
\oddsidemargin=0in
\textwidth=6.5in
\textheight=9.0in
\headsep=0.25in 
\numberwithin{equation}{section}
\linespread{1.1} 
\setlength\parindent{0pt} % Removes all indentation from paragraphs
\setlength\answerskip{-.075in}
\setlength\answerlinelength{.5in}
\newcommand{\myitem}{\part[-]}




% HEADER AND FOOTER
\pagestyle{headandfoot}
\runningheadrule
\runningfootrule
\firstpageheader{Kutztown University}{  }{MATH 181 -- Calculus I}
\firstpageheadrule
\runningheader{Desmos Lab 4}{ }{MATH 181}
\firstpagefooter{}{}{} 
\runningfooter{}{}{Page \thepage\ of \numpages}


\newcommand{\assignmentname}{Desmos Lab IV} % Assignment title
\newcommand{\coursetitle}{MATH 181 -- Calculus I}
\newcommand{\duedate}{Due: Friday, December 5} % Class/lecture time
\newcommand{\readsections}{Read: Section 5.3 and 5.5} % Teacher/lecturer
\newcommand{\semester}{Fall 2025}
\title{
\vspace{2in}
\textmd{\textbf{\course:\ \assignmentname}}\\
\normalsize\vspace{0.1in}{\duedate}\\
\vspace{0.1in}\large{\textit{\instructor}}
\vspace{3in}
}
\author{\textbf{Author: \authorname}}
\date{Last Updated: \today}

\pointsinrightmargin
\bracketedpoints


%\printanswers % Unomment this line to Print Solutions 

\begin{document}


%%%%%%%%%%%%%%%%%%%%%%%%%%%%%%%%%%%%%%%%%%%%%%%%%%%%%
% TITLES AND SUCH
%%%%%%%%%%%%%%%%%%%%%%%%%%%%%%%%%%%%%%%%%%%%%%%%%%%%%
\hspace{1in}\\ \vspace{.15in}


\hbox to \textwidth{Name: \hrulefill\hrulefill }

\begin{center}
\normalsize
 \fbox{\fbox{\parbox{5.5in}{

        {\bf Instructions:}  All solutions should be prepared carefully and neatly.  Any questions that ask you to explain or justify should have answers written in complete sentences.  All solution sets shall be completed on this packet and handed in at the beginning of class on the due date below.  Careless presentation (e.g.~bad handwriting, pen scribbles, doodles, wasted space, etc) will result in a deduction of points at my discretion.  If the question asks to submit a graph or plot, be sure to print the plot and include it in your final lab report.  Use a PENCIL and if you make a mistake, use an eraser.  Each student should prepare their own solutions that demonstrate their own comprehension of the material.  Submitted work that does not demonstrate clearly the process by which one arrived at the answer will not receive credit of any kind.  Academic dishonesty will not be tolerated. }}}
\end{center}

\vspace{.5cm}

\ifprintanswers
	\begin{center}
		\Huge\textsc{\assignmentname}\\
		\Large{\textsc{\coursetitle}}\\[4pt]
		\Large\textsc{\color{red}Solution Key
		}
	\end{center}
\else
	\begin{center}
		\Huge\textsc{
		\assignmentname}\\[4pt]
                 \Large{\textsc{\coursetitle}\\[10pt]}
                 \Large{\textsc{\duedate}\\[8pt]}
                 \Large{\textsc{\readsections}\\[8pt]
                 }
	\end{center}
\fi

\Large
\begin{center}
\textsc{\textcolor{red}{Solutions}}
\end{center}


%%%%%%%%%%%%%%%%%%%%%%%%%%%%%%%%%%%%%%%%%%%%%%%%%%%%%
%%%%%%%%%%%%%%%%%%%%%%%%%%%%%%%%%%%%%%%%%%%%%%%%%%%%%


\clearpage
%%%%%%%%%%%%%%%%%%%%%%%%%%%%%%%%%%%%%%%%%%%%%%%%%%%%%
% Blank Page
%%%%%%%%%%%%%%%%%%%%%%%%%%%%%%%%%%%%%%%%%%%%%%%%%%%%%
\normalsize
\
\vspace{.5cm}

\huge
\begin{center}
\begin{tabular}{|c|c|c|}
\hline
Problem  & Available  &   Your  \\
Number & Points &    Points  \\	\hline\hline

 $1$  & 10  &  \textcolor{red}{10} \\	\hline
 $2$  & 10  &   \textcolor{red}{10}  \\	\hline

Total & 20  &  \textcolor{red}{20}\\ \hline
\end{tabular}
\end{center}

\pagebreak

\normalsize
\begin{questions}


%%%%%%%%%%%%%%%%%%%%%%%%%%%%%%%%%%%%%%%%%%%%%%%%%%%%
%Part 1
%%%%%%%%%%%%%%%%%%%%%%%%%%%%%%%%%%%%%%%%%%%%%%%%%%%%
\question \underline{Definite Integrals:} In this part, we'll use Desmos to compute definite integrals and validate the constant multiple property of integrals.

\begin{parts}
	\part Open a new graph in Desmos and save it as ``Definite Integrals."
	\part Define the following function on the first line in Desmos: $\displaystyle f(x) = 4x^3-2x +2$.  Define the domain of the function from $-1$ to $1$ by typing $\{-1\leq x \leq 1\}$ right next to the function.  Graph the function on the window $-2\leq x \leq 2$ and $-0.5 \leq y\leq 5$.
	\part On the next line, type $0\leq y\leq f(x)\, \{-1\leq x \leq 1\}$.  What does Desmos do?  Explain briefly below. 
	
	\textcolor{red}{It shades in the area under the curve but above the $x$-axis. }
	
	\part Compute the integral $\displaystyle \int_{-1}^1 4x^3 - 2x +2\, dx$ by hand in the space below using the Fundamental Theorem of Calculus.
	
	\textcolor{red}{The integral using the Fundamental Theorem of Calculus is 
	\begin{align*}
		\int_{-1}^1 4x^3 - 2x + 2\, dx & = \left. \frac{4}{4}x^4 - \frac{2}{2}x^2 + 2x\right|_{-1}^1 \\ 
		& = \left. x^4 - x^2 + 2x\right|_{-1}^1 \\ 
		& = \big(1^4 - 1^2 + 2(1)\big) - \big((-1)^4 - (-1)^2 + 2(-1)\big) \\ 
		& = 1 - 1 + 2 - 1 + 1 + 2 \\ 
		& = 4 \end{align*}}
	\part Now, on a new line, use Desmos to check your answer by typing in the integral exactly how it looks above. 
	\part Graph the area below $y = 2f(x)$ by typing $0\leq y\leq 2f(x)\, \{-1\leq x \leq 1\}$.  Have Desmos calculate this area.  Is the answer what you expect?  Explain briefly.
	
	\textcolor{red}{The area below the new function $y = 2f(x)$ is exactly double the area of the original, which is exactly what we expected because we can take the constant through the integral.  Taking the 2 out allows us to simply find the area of the original region, which was 4, and double it.  }
	
		\part When you submit this assignment on D2L, share the \textbf{link} to this graph

			\textcolor{red}{\texttt{https://www.desmos.com/calculator/zxppicgypo}}
	

\end{parts}

\pagebreak



%%%%%%%%%%%%%%%%%%%%%%%%%%%%%%%%%%%%%%%%%%%%%%%%%%%%
%Part 2
%%%%%%%%%%%%%%%%%%%%%%%%%%%%%%%%%%%%%%%%%%%%%%%%%%%%
\question \underline{Symmetric Functions and Area:} In this part, we'll use Desmos to explore how areas bounded by odd and even functions result in a somewhat easier calculation. 
\begin{parts}
	\part Open a new graph in Desmos and save it as ``Symmetric Functions."
	\part Create a folder called ``Even Functions."  Define the following function on the first line in the folder: $\displaystyle f_e(x) = -3x^4+4x^2+2$.  Define the domain of the function from $-a$ to $a$ by typing $\{-a\leq x \leq a\}$ right next to the function.  Define a slider for $a$ and let $a$ range from 0 to 5. 
	\part On the next two lines in the folder, type $0\leq y\leq f_e(x)\, \{-a\leq x \leq a\}$ and $f_e(x)\leq y\leq 0\, \{-a\leq x \leq a\}$. Push play on $a$.  What does Desmos do?  Explain briefly below. 
	
	\textcolor{red}{It graphs the function $y = f_e(x)$ and expands the left and right endpoints out, while at the same time, shading in the region under the curve above the $x$-axis and above the curve below the $x$-axis.  Essentially, it is showing the definite integral of the function over $-a$ to $a$.  }
	
	\part Outside the folder, have Desmos calculate the following integral:  $\displaystyle \int_{-a}^a f_e(x), dx$.  Then on a new line, have it calculate $\displaystyle \int_{0}^a f_e(x), dx$.  What is the relationship between these two areas?  Explain.

	\textcolor{red}{For any value of $a$, the area $\displaystyle \int_{-a}^a f_e(x)\, dx$ is exactly double the area $\displaystyle \int_0^a f_e(x)\, dx$.  This is because the graph of $y = f_e(x)$ is symmetric about the $y$-axis.  The area contained by the graph from $0$ to $a$ is essentially identical to the area contained between $-a$ and $0$. Therefore, we need only multiply the area from $0$ to $a$ by 2 to get the full area.    }

	\part Toggle off the graph of the even function by clicking the folder.  Now, create a new folder called ``Odd Functions."  Define an odd function of your choice in this folder and call it $f_o(x)$.  Plot the area between $-a$ and $a$, just like you did for the even function.  Now, have Desmos calculate the area $\displaystyle \int_{-a}^a f_o(x), dx$.  Is the area independent of $a$?  Why?  Explain briefly.  

\textcolor{red}{The area contained from $-a$ to $a$ for an odd function like $f_o(x)$ will always be zero.  Yes, the area is independent of the value of $a$.  This is because the function values on the right side of the origin are the opposite in sign to those on the left side of the origin.  Therefore, when we compute the area bounded by $f_o(x)$ from $-a$ to $0$ and add it to the area from $0$ to $a$, the two areas are equal in magnitude but opposite in sign.  Therefore, the full area from $-a$ to $a$ will always total $0$. }
	
	
	\part When you submit this assignment on D2L, share the \textbf{link} to this graph

				\textcolor{red}{\texttt{https://www.desmos.com/calculator/hxn9skmyth}}


\end{parts}






\end{questions}

\end{document} 














